\scanpage{75}
\mspage{66}
% archive sheet number: 75

la suivante :
\[
\begin{aligned}
\operatorname{cosupp}X&=\inserted{\{Y\in M\mid\widetilde{Y}\subset}\operatorname{cosupp}^{\circ}(\widetilde{X})\inserted{\}}\\
\operatorname{cosupp}A&=\inserted{\{Y\in M\mid\widetilde{Y}\subset}\operatorname{cosupp}^{\circ}(\widetilde{A})\inserted{\}}
\end{aligned}
\]
(où $\widetilde{A}\overset{\mathrm{def}}{=}\bigcup_{X\in A}\widetilde{X}$ (\deleted{\unsure{saturée de}}\inserted{plus petite} partie fermée de $M,\leq$ contenant $A$)

\deleted{Donc $\operatorname{cosupp}A=\operatorname{cosupp}^{\circ}(A)$ si $A=\widetilde{A}$ i.e. $A$ fermé dans $M$ pour $\leq$.}
% margin note written slanted
\margin{\uline{idiot}}

% Two long oblique cancellation strokes cross the following block.
\deleted{\unsure{En.} Ainsi, les supports sont les supports \unsure{internes} des parties de $M$ \uline{fermées} pour $\leq$. En effet, pour une telle partie, prendre $\operatorname{cosupp}^{\circ}$ ou $\operatorname{cosupp}$ revient au $=$ et le résultat est $\leq$-fermé, et \unsure{le réappliquant}, on trouve l'assertion.}

\deleted{Le traitement précédent des supports était vicié par le fait, qu'il ne travaillait \inserted{seulement} \deleted{\unsure{au démarrage}} qu'avec les supports des parties fermées.} Les $\operatorname{supp}$ et $\operatorname{cosupp}.$ s'expriment, on l'a vu, en termes des $\operatorname{supp}^{\circ}$ et $\operatorname{cosupp}^{\circ}$, mais l'inverse n'est pas vrai.

\uline{Exemple instructif} : \unsure{magma} trivial\inserted{\unsure{ou n'importe,}}. La relation $\mid\circ\mid$ est la relation $X\ne Y$. Donc les supports sont les parties quelconques de $M$. Je peux visualiser une telle partie $A$, comme représentant la réunion des ``\uline{cellules ouvertes}'' $X^{\circ}$, pour $X\in A$. Si on
